\subsection{W-21 (VWM)}

Techniky pro vyhledávání textových, webových a multimediálních dokumentů: modely, algoritmy, aplikace. Optimalizace webových stránek pro vyhledávače.

\subsubsection*{Základní pojmy}

Crawling - stahování obsahu stránek
indexing - úprava dat do podoby, ve které se dobře vyhledává
Searching - hledání na základě dotazu

dokument - entita obsahující text
kolekce - množina dokumentů

Vyhledávání v plaintextu - obří množství dat, které je potřeba prohledat - je nutné využít index.

\subsubsection*{Optimalizace indexování}

- odstranění stopwords - často se opakující slova (the, is, a), ale i čísla nebo jména
- stemming - zjednodušení slov (např. "running" na "run")

\subsubsection*{Metriky vyhledávání}

- effectivness - míra uspokojení uživatele s výsledky
- precision - pravděpodobnost, že dokument ve výsledku je relevantní
- recall - pravděpodobnost, že relevantní dokument je ve výsledku

precision klesá s vyšším recallem.

\subsubsection*{Boolean retrieval model}

- založený na logice "dokument obsahuje nebo neobsahuje slovo"
- dotazovatelný "jazyk" pomocí operátorů AND, OR, NOT - Caesar AND salad
- implementováno pomocí invertovaného indexu v podobě výraz-seřazený seznam dokumentů ve kterých se vyskytuje

lze optimalizovat:
- ke každému výrazu si pamatujeme počet výskytů
- časté výrazy se neindexují, protože by měly dlouhé seznamy

výhody: jednoduché na implementaci, profesionálové (právníci) často využívají, protože ví přesně co hledají
nevýhody: vysoká přesnost dotazu znamená nízkou efektivitu (precision a recall), dokumenty nejsou seřazené podle relevance - nemají hodnocení

existuje rozšířený boolean model, který využívá score výrazů - kombinuje boolean retrieval model a vektorové vyhledávání (vzdálenosti)

\subsubsection*{Vector retrieval model}

- dokumenty a dotazy jsou reprezentovány jako vektory
- euklidovská vzdálenost se nepoužívá, lepší je cosine
- jedna z nejlepších metod vážení je TF-IDF - term frequency - inverse document frequency (čím častěji se slovo vyskytuje v dokumentu, tím je důležitější, čím méně se vyskytuje v kolekci, tím je důležitější)
- výsledkem je matice dokumentů a výrazů, kde každý dokument je reprezentován jako vektor výrazů

\subsubsection*{Latent semantic indexing}

- rozšíření vektorového modelu
- dokument je reprezentován jako vektor, ale místo výrazů se používají latentní koncepty
- latentní koncepty jsou získány pomocí SVD (singular value decomposition)
- SVD je metoda pro redukci dimenzí, která se používá k nalezení latentních konceptů (převede "sparse" vektory na "dense" vektory "vybráním" nejdůležitějších dimenzí)
- výsledná matice dokumentů a konceptů je dense a menší než původní matice 
- náročný výpočet, ale dobré výsledky pokud hledáme koncepty / významy

\subsubsection*{Word2Vec}

- machine learning metoda na získávání vektorových reprezentací slov (případně genů, kódu, a dalších druhů dat)
- využívá dvě neuronové sítě
- rychlejší než LSI, má i lepší výsledky (window contexts, not entire documents)
- vec("king") - vec("man") + vec("woman") = vec("queen")
